\documentclass[output=paper,colorlinks,citecolor=brown]{langscibook}
\author{Marine Vuillermet\orcid{}\affiliation{University of Zürich}}
\title{A fieldwork questionnaire on apprehensionals}
\abstract{This questionnaire is intended as a guide to fieldworkers for systematically surveying the constructions used in the apprehensional functions known so far, namely apprehensive, precautioning, `fear' complementation and timitive. It builds on the parameters of variation attested so far crosslinguistically, at the semantic, syntactic and pragmatic levels, and comes in the form of brief explanations of the variation targeted, illustrated by prompts in English to be translated into the language of study. A systematic investigation of the fine-grained semantics of the different functions of apprehensionals may uncover subfunctions that have not been clearly identified yet, and/or offer an aid in establishing whether or not these expressions are dedicated apprehensionals.}


\IfFileExists{../localcommands.tex}{
   \addbibresource{../localbibliography.bib}
   \usepackage{tabularx,multicol}
\usepackage{url}
\urlstyle{same}

\usepackage{langsci-optional}
\usepackage{langsci-lgr}
\usepackage{langsci-gb4e}

\babelprovide[import]{hebrew}

\usepackage{paralist} %%may not be needed, currently used in questionnaire


%Siena added packages for introduction; not sure they are needed
\usepackage[normalem]{ulem}
\useunder{\uline}{\ul}{}

%from Treis:
\usepackage{listings}
\lstset{basicstyle=\ttfamily,tabsize=2,breaklines=true}

%%from Francez
% % % \usepackage{cjhebrew}
%\usepackage{tipa}  Also included by D&A, need to check with LSP whether to change it to language science tipa
\usepackage{leipzig}
%%from Dabkowski&AnderBois:

\usepackage[shortlabels]{enumitem} 
\usepackage{centernot}
\usepackage{arydshln}
\usepackage{newunicodechar}
\usepackage{csquotes}
%\let\sups\relax \usepackage{tipa} commented out on 7/3/25 Martina
% \usepackage[all]{nowidow}
\usepackage{listings} \lstset{basicstyle=\ttfamily}
\usepackage{multirow}
\usepackage{rotating}
\usepackage{pdflscape}
\usepackage{xspace}
%\usepackage{tabularx}
%\usepackage{multicol}
\usepackage{supertabular}
\usepackage{hhline}
\usepackage{bbding} %this package is for the checkmarks in Table 6 in Browne et al.

%\usepackage{qtree}
%\usepackage{tree-dvips}

\usepackage{array}
\usetikzlibrary{calc,tikzmark,matrix}
\usepgflibrary{shadings}
\usepackage{pifont}
%from Fedotov:

\usepackage{langsci-textipa}
%\usepackage{multirow}


%\usepackage{xeCJK}

%%%%%%%%%%%%%%%%%%%%%%%%%%%%%%%%%%%%%%%%%%%%%%%%%%%%
%%% EXAMPLES %%%%%%%%%%%%%%%%%%%%%%%%%%%%%%%%%%%%%%%
%%%%%%%%%%%%%%%%%%%%%%%%%%%%%%%%%%%%%%%%%%%%%%%%%%%% 

% if you want the source line of examples to be in
% italics, uncomment the following line
\renewcommand{\exfont}{\itshape}
% to add additional information to the right of
% examples, uncomment the following line

%\usepackage{langsci/styles/jambox}


%%% FOOTNOTE EXAMPLE NUMBERING
% \makeatletter
% \patchcmd{\@footnotetext}{\setcounter{fnx}{0}}{\renewcommand{\thexnumi}{\roman{xnumi}}}{}{}
% \apptocmd{\@footnotetext}{
%     \@noftnotetrue
%     \renewcommand{\thexnumi}{\arabic{xnumi}}%
% }{}{}
%
% \makeatother
\usepackage[linguistics,edges]{forest}
\usepackage{hyphenat}
\usepackage{langsci-branding}

   \makeatletter
\let\thetitle\@title
\let\theauthor\@author
\makeatother

\newcommand{\togglepaper}[1][0]{
   \bibliography{../localbibliography}
   \papernote{\scriptsize\normalfont
     \theauthor.
     \titleTemp.
    To appear in:
    Martina Faller,  Marine Vuillermet \&  Eva Schultze-Berndt (eds.). Apprehensional Constructions in a cross-linguistic perspective. Berlin: Language Science Press. [preliminary page numbering]
   }
   \pagenumbering{roman}
   \setcounter{chapter}{#1}
   \addtocounter{chapter}{-1}
}

\newbool{bookcompile}
\booltrue{bookcompile}
\newcommand{\bookorchapter}[2]{\ifbool{bookcompile}{#1}{#2}}


\newfontfamily\FedotovBoldEmptySetFont{LibertinusSerif-Italic.otf}[AutoFakeBold=2.5]
\newcommand{\FedotovBoldEmpty}{{\FedotovBoldEmptySetFont\bfseries ∅}}

%From Dabkowski&AnderBois


\let\sc\scshape
\let\rm\rmfamily
%\let\it\itshape
\let\tt\ttfamily
\let\nr\normalfont

\newcommand{\wsc}[1]{\textls[80]{\scshape#1}}

\let\textai\textit
% \let\textlc\spacedlowsmallcaps
% \let\textac\spacedallcaps
\let\textlc\textsc
\let\textac\textsc


%%%%%%%%%%%%%%%%%%%%%%%%%%%%%%%%%%%%%%%%%%%%%%%%%%%%
%%% FLAG %%%%%%%%%%%%%%%%%%%%%%%%%%%%%%%%%%%%%%%%%%%
%%%%%%%%%%%%%%%%%%%%%%%%%%%%%%%%%%%%%%%%%%%%%%%%%%%%

% \newcommand{\scott}  [1]{\textcolor{violet}{#1}}
% \newcommand{\maks}   [1]{\textcolor{blue}{#1}}
% % \newcommand{\data}   [1]{\textcolor{teal}{#1}}
% \newcommand{\new}   [1]{#1}
% \newcommand{\newnew}   [1]{\textcolor{teal}{#1}}
% \newcommand{\rev}   [1]{\textcolor{magenta}{#1}}



%%%%%%%%%%%%%%%%%%%%%%%%%%%%%%%%%%%%%%%%%%%%%%%%%%%%
%%% TABLES %%%%%%%%%%%%%%%%%%%%%%%%%%%%%%%%%%%%%%%%%
%%%%%%%%%%%%%%%%%%%%%%%%%%%%%%%%%%%%%%%%%%%%%%%%%%%%

%\newcolumntype{Y}{>{\arraybackslash}p{25.3543464286pt}} % length of L

% \newcolumntype{K}{>{\arraybackslash}p{24.25pt}}
% \newcolumntype{V}{>{\arraybackslash}p{(\textwidth/9-24pt)/2}}
%
\let\ch\relax \newcommand{\ch}[2]{\multicolumn{#1}{c}{\sc#2}}
\let\sx\relax \newcommand{\sx}[2]{\multicolumn{#1}{l}{\emph{-#2}}} % suffix
\let\su\relax \newcommand{\su}{\sx{1}} % suffix unicolumnal
\let\sd\relax \newcommand{\sd}{\sx{2}} % suffix dicolumnal
\let\cx\relax \newcommand{\cx}[2]{\multicolumn{#1}{l}{\emph{=#2}}} % clitic
\let\cu\relax \newcommand{\cu}{\cx{1}} % clitic unicolumnal
\let\cd\relax \newcommand{\cd}{\cx{2}} % suffix dicolumnal
\let\gx\relax \newcommand{\gx}[2]{\multicolumn{#1}{l}{\sc #2}} % gloss
\let\gd\relax \newcommand{\gd}{\gx{2}} % gloss dicolumnal
\let\gr\relax \newcommand{\gr}[1]{\multicolumn{2}{c}{\multirow{2}{*}{\textcolor{gray}{\sc#1}}}}
\let\db\relax \newcommand{\db}[1]{\multicolumn{2}{c}{\multirow{2}{*}{\sc#1}}}



%%%%%%%%%%%%%%%%%%%%%%%%%%%%%%%%%%%%%%%%%%%%%%%%%%%%
%%% EXAMPLES %%%%%%%%%%%%%%%%%%%%%%%%%%%%%%%%%%%%%%%
%%%%%%%%%%%%%%%%%%%%%%%%%%%%%%%%%%%%%%%%%%%%%%%%%%%%

%%% AFFIX BOUNDARY %%%%%%
% \DeclareRobustCommand{\longmn}{\textminus}
% \DeclareRobustCommand{\mn}{-}

%%% CLITIC BOUNDARY %%%%%%
% \DeclareRobustCommand{\longeq}{=}
% \makeatletter\DeclareRobustCommand{\eq}{%
%     \settowidth{\@tempdima}{-}% width of a hyphen
%     \resizebox{\@tempdima}{\height}{=}%
% }\makeatother

%%% FUZZY CLITIC %%%%%%
% \DeclareRobustCommand{\longfc}{%
%     $\overset{\text{\smash{}}}{\text{$\approx$}}$%
% }
% \makeatletter\DeclareRobustCommand{\fc}{%
%     \settowidth{\@tempdima}{-}% width of a hyphen
%     \resizebox{\@tempdima}{\height}{%
%         $\overset{\text{\smash{}}}{\text{$\approx$}}$%
%     }%
% }\makeatother

%%% REDUPLICATION %%%%%%%
% \DeclareRobustCommand{\longtl}{$\sim$}
% \makeatletter\DeclareRobustCommand{\tl}{%
%     \settowidth{\@tempdima}{-}% width of a hyphen
%     \resizebox{\@tempdima}{\height}{$\sim$}%
% }\makeatother

% \newcommand{\mn}          {\textminus}
\newcommand{\src}    [1]  {\jambox[0pt]{\rm(#1)}}
% \newcommand{\ai}     [2]  {\emph{#1} `\textsc{#2}#3'\xspace}
\newcommand{\ai}     [2]  {\emph{#1} `\textsc{#2}'\xspace}
\newcommand{\sane}   [1][]{\ai{=sa'ne}{appr}{#1}}
\newcommand{\srcn}   [1][]{\ai{kûintsû}{srcn}{#1}}
% \newcommand{\mbekane}[1][]{\ai{\eq mbe kañe}{neg.adv aux.inf}{#1}}

% \let\sc\relax \newcommand{\sc}{\normalfont\scshape}
% \let\rm\relax \newcommand{\rm}{\normalfont\rmfamily}
% \let\it\relax \newcommand{\it}{\normalfont\itfamily}

\newcommand{\lb}{{\rm[}}
\newcommand{\rb}{{\rm]}}

\newcommand{\bad}{\makebox[0pt][r]{*\,}}
\newcommand{\infel}{\makebox[0pt][r]{\#\,}}


% \newcommand{\lsqz}[1]{#1}
% \newcommand{\lsqz}[1]{\makebox[0pt][l]{#1}}

% \newcommand{\lsqu}[1]{\makebox[0pt][l]{#1}}
% \newcommand{\rsqu}[1]{\makebox[0pt][r]{#1}}
% \newcommand{\astr}{\makebox[0pt][r]{{\nr*}}}
% \newcommand{\hash}{\makebox[0pt][r]{{\nr\#}}}
% \newcommand{\qstn}{\makebox[0pt][r]{{\nr?}}}
% \newcommand{\bad}{\makebox[0pt][r]{*}}
% \newcommand{\unhappy}{\makebox[0pt][r]{\#}}


%%%%%%%%%%%%%%%%%%%%%%%%%%%%%%%%%%%%%%%%%%%%%%%%%%%%
%%% UNICODE CHARACTERS  %%%%%%%%%%%%%%%%%%%%%%%%%%%%
%%%%%%%%%%%%%%%%%%%%%%%%%%%%%%%%%%%%%%%%%%%%%%%%%%%%

\newunicodechar{⟨}{⟨}
\newunicodechar{⟩}{⟩}

\newcommand{\infix}{⟨}
\newcommand{\outfix}{⟩}



%%%%%%%%%%%%%%%%%%%%%%%%%%%%%%%%%%%%%%%%%%%%%%%%%%%%
%%% IPA %%%%%%%%%%%%%%%%%%%%%%%%%%%%%%%%%%%%%%%%%%%%
%%%%%%%%%%%%%%%%%%%%%%%%%%%%%%%%%%%%%%%%%%%%%%%%%%%%

% \newcommand{\ipa}{\textipa}

\newcommand{\sub}{\textsubscript}
% \let\super\relax \newcommand{\super}{\textsuperscript}

\newcommand{\ts}{\texttslig}
% \let\dz\relax \newcommand{\dz}{\textdzlig}
\newcommand{\tS}{\textteshlig}
\newcommand{\dZ}{\textdyoghlig}
\newcommand{\cj}{\textctc}
\newcommand{\sj}{\textctc}
\newcommand{\zj}{\textctz}
\newcommand{\tcj}{\texttctclig}
\newcommand{\tsj}{\texttctclig}
\newcommand{\dzj}{\textdctzlig}
\newcommand{\gj}{\textbardotlessj}
% \let\nj\relax \newcommand{\nj}{\textltailn}
% \newcommand{\W}{\textturnmrleg}
\newcommand{\wl}{\textturnmrleg}
% \newcommand{\1}{\ipabar{\~\i}{.5ex}{1.1}{}{}}
\newcommand{\dotlessibar}{\ipabar{\i}{.5ex}{1.1}{}{}}
\newcommand{\darkl}{\textltilde}
% \newcommand{\n}{\textsubarch}
% \newcommand{\x}{\textovercross}
\newcommand{\lowered}{\textlowering}
\newcommand{\raised}{\textraising}
\newcommand{\retracted}{\textsubbar}
\newcommand{\unreleased}{\textcorner}
% \newcommand{\syllabic}{\s}
\newcommand{\implosivet}{\texthtt}
\newcommand{\implosived}{\texthtd}
\newcommand{\implosivep}{\texthtp}
\newcommand{\implosiveb}{\texthtb}
\newcommand{\implosivek}{\texthtk}
\newcommand{\implosiveg}{\texthtg}



%%%%%%%%%%%%%%%%%%%%%%%%%%%%%%%%%%%%%%%%%%%%%%%%%%%%
%%% OTHER %%%%%%%%%%%%%%%%%%%%%%%%%%%%%%
%%%%%%%%%%%%%%%%%%%%%%%%%%%%%%%%%%%%%%%%%%%%%%%%%%%%

\newcommand{\ie}{i.\,e.\xspace}
\newcommand{\Ie}{I.\,e.\xspace}
\newcommand{\eg}{e.\,g.\xspace}
\newcommand{\Eg}{E.\,g.\xspace} 




\newcommand{\francoissource}[1]{\hfill\textnormal{{\footnotesize #1} }}
\newcommand{\E}{Ɛ\kern-1.5pt\raisebox{1pt}{̏}\kern1.5pt}
\newcommand{\ù}{{ṵ}\kern-2pt{̀}\kern2pt}
\newcommand{\ȕ}{{ṵ}\kern-2pt{̏}\kern2pt}

\newcommand{\OO}{ɔ\kern-1pt\raisebox{-1pt}{̰}\kern1pt}
\newcommand{\Ǒ}{{{ɔ̌\kern-1pt\raisebox{-.5pt}{̰}\kern1pt}}}


\newcommand{\à}{{a̰}\kern-2pt{̀}\kern2pt}
\newcommand{\ilit}[1]{#1\il{#1}}

   %% hyphenation points for line breaks
%% Normally, automatic hyphenation in LaTeX is very good
%% If a word is mis-hyphenated, add it to this file
%%
%% add information to TeX file before \begin{document} with:
%% %% hyphenation points for line breaks
%% Normally, automatic hyphenation in LaTeX is very good
%% If a word is mis-hyphenated, add it to this file
%%
%% add information to TeX file before \begin{document} with:
%% %% hyphenation points for line breaks
%% Normally, automatic hyphenation in LaTeX is very good
%% If a word is mis-hyphenated, add it to this file
%%
%% add information to TeX file before \begin{document} with:
%% \include{localhyphenation}

\hyphenation{
par-a-digm
com-ple-men-ta-tion
anaph-o-ra
Dor-drecht
Cu-shi-tic
be-ne-dic-tive
pa-ra-digms
Ethi-o-pi-an
hoch-land-ost-ku-schi-ti-schen 
Duu-raa-me
Pa-ken-dorf
Lich-ten-berk
affri-ca-te
affri-ca-tes
com-ple-ments
Eng-lish
}




\hyphenation{
par-a-digm
com-ple-men-ta-tion
anaph-o-ra
Dor-drecht
Cu-shi-tic
be-ne-dic-tive
pa-ra-digms
Ethi-o-pi-an
hoch-land-ost-ku-schi-ti-schen 
Duu-raa-me
Pa-ken-dorf
Lich-ten-berk
affri-ca-te
affri-ca-tes
com-ple-ments
Eng-lish
}




\hyphenation{
par-a-digm
com-ple-men-ta-tion
anaph-o-ra
Dor-drecht
Cu-shi-tic
be-ne-dic-tive
pa-ra-digms
Ethi-o-pi-an
hoch-land-ost-ku-schi-ti-schen 
Duu-raa-me
Pa-ken-dorf
Lich-ten-berk
affri-ca-te
affri-ca-tes
com-ple-ments
Eng-lish
}



   \boolfalse{bookcompile}
   \togglepaper[23]%%chapternumber
}{}


\begin{document}
\maketitle

\section*{Introduction}
%\label{sec:intro}

Apprehensionals is a cover term that refers to any grammatical construction expressing the (high) \textit{possibility} of an \textit{undesirable} event, or expressing the undesirability of an entity. This questionnaire considers in turn the following four functions: \textsc{Apprehensive}, \textsc{Precautioning}, \textsc{`fear' complementation}, and \textsc{Timitive}. Their similarities and differences are summarized in the table below, and clarified in the introduction of each (sub)section.


\begin{table}
\small
\caption{Four functions of apprehensionals}
\label{tab:functions}
\begin{tabularx}{\textwidth}{QlQ>{\raggedright}p{18mm}l}
\lsptoprule
{} & {Scope} & {Obligatory presence of} & {Typical\newline function} & {Section} \\
\midrule
\textbf{Apprehensive}\newline \textit{\textbf{watch out} you \textbf{might} fall} & clause & NA & warning, threat & \sectref{sec:apprehensive} \\
\tablevspace
\textbf{Precautioning}\newline \textit{tie your laces \textbf{lest} you fall} & clause & a preemptive clause to which the undesirable event is associated & justification of a preemptive action & \sectref{sec:precautioning} \\
\tablevspace
\textbf{\mbox{`fear' complementation}}\newline \textit{\textbf{I fear lest} you fall} & clause & the undesirable event is the non-subject complement of a transitive predicate of fear & account of the feeling of fear & \sectref{subsec:prefear} \\
\tablevspace
\textbf{Timitive}\newline \textit{she ran away \textbf{for fear}} \textit{of the dog} & NP & a main clause encoding the reaction wrt the undesirable entity & justification of a reaction & \sectref{sec:timitive} \\
\lspbottomrule
\end{tabularx}%
\end{table}


Languages differ a lot in the number of apprehensional constructions they have -- up to three in Ese Ejja (Pano-Tacanan; Bolivia \& Peru; \cite{vuillermet2018a}) -- as well as in the number of functions (synchronically) available to their apprehensional construction(s) -- up to three functions for a single apprehensional marker in Marrithiyel (Western Daly; Australia; \cite{green1989}).


\ea \label{ex:v1} \textnormal{Ese Ejja and its three apprehensionals (Apprehensive, Precautioning, Timitive)}
\ea \label{ex:v1a} \textnormal{\textsc{Apprehensive}} \\
\gll 'Biya 'biya 'biya 'biya! Kekwa-ka-\textbf{chana} miya! \\
bee bee bee bee pierce-\textsc{3a}-\textbf{\textsc{appr}} \textsc{2sg.abs} \\
\glt `Bee, bee, bee, bee! \textbf{Watch out} it \textbf{might} sting you!'\footnote{The combination of \textit{watch out} or \textit{be careful} with the modal \textit{might} is probably the best proxy in English for an apprehensive sentence. The modal \textit{might} expresses potentiality only, without further specifying whether it is positively or negatively evaluated: \textit{he might come} might be positive, e.g.\ in \textit{that he might come makes me so happy!} The round brackets signal that equivalents of these two phrases are not necessary in languages with dedicated apprehensive markers, although they might often co-occur.} \\
\ex \label{ex:v1b} \textnormal{\textsc{Precautioning}} \\
\gll Owaya ekowijji shijja-ka-ani {\ob}\textbf{e}-jja-saja-ki \textbf{kwajejje}{\cb}.\\
\textsc{3erg} rifle clean-\textsc{3a-prs} \textsc{\textbf{prec}-mid-}block-\textsc{mid} \textbf{\textsc{prec}} \\
\glt `He cleans his rifle [\textbf{lest} it get blocked].'
\newpage
\ex \label{ex:v1c} \textnormal{\textsc{Timitive}} \\
\gll Iñawewa kwaji\textasciitilde kwaji-ani 'biya=\textbf{yajjajo}. \\
dog run\textasciitilde\textsc{rdp-prs} bee=\textsc{\textbf{tim}} \\
\glt `The dog is running \textbf{for fear of} the bees.'
\z
\z



\ea \label{ex:v2} \textnormal{Marrithiyel and its multifunctional apprehensional marker \citep[80, 166, 58, MV's emphasis]{green1989}}
\ea \label{ex:v2a} \textnormal{\textsc{Apprehensive} function} \\
\gll Gu-n-ning-pirr-Ø-\textbf{fang}. \\
\textsc{3s.real}-go-\textsc{1sg.o.nsg.nirr.s}-leave-\textsc{3pl-\textbf{apprl}} \\
\glt `({\textbf{I am afraid}}) they might leave me.' \\
\ex \label{ex:v2b} \textnormal{\textsc{Precautioning} function }\\
\gll \textnormal{(\ldots)} {\ob}atjan ambi gu-iwinj-sjang-Ø-\textbf{fang}{\cb}.\\
{} dog \textsc{neg} \textsc{3nsg.s.real-3nsg-}listen-\textsc{3pl-\textbf{apprl}} \\
\glt `(He turned the generator off), [\textbf{lest} they not hear the dogs].'
\ex \label{ex:v2c} \textnormal{\textsc{Timitive} function }\\
\gll Wanthi ngun-wa ngaka yigin-\textbf{fang}. \\
behind \textsc{1sg.sbj.irr}-go\textsc{.fut} sister \textsc{1sbj.poss-\textbf{apprl}} \\
\glt `I will come afterwards, \textbf{to avoid my sister}.'
\z
\z

The questionnaire builds on a previous manuscript \citep{VuillermetM2018questionnaire} but includes substantial revisions and additions.  It is intended as a guide to fieldworkers for surveying the constructions used in the target language in the apprehensional functions listed above, and as an aid in establishing whether or not these expressions are dedicated apprehensionals. It  systematically tests for the parameters of variation attested so far crosslinguistically, at the semantic, syntactic and pragmatic levels.

It consists of prompts in English to be translated into the language of study. The targeted grammatical morpheme is paraphrased in English if no equivalent exists (as for the timitive examples in \ref{ex:v1c} and \ref{ex:v2c}). Prompts are preceded by an explanation of the targeted parameter(s) and, sometimes, by a reference to a specific language illustrating the parameter(s) and/or suggestions of possible contexts to facilitate the elicitation task. Ideally, the researcher should suggest the form in the target language, and ask the speaker to:

\begin{enumerate}
    \item Give a context when the sentence could be uttered -- contexts should be suggested if the speaker cannot find any;
    \item (Where applicable) translate it back into a national or regional language and/or the language of work.
    
\end{enumerate}

\noindent Ideally again, the sentences should be checked with different speakers.

To avoid cultural or situational incongruities, several prompts are offered, among which the researcher might choose. They may also create their own prompt with the same variation target in mind to adapt the questions to the consultant, language, and/or situation. As an example, restrictions on the grammatical persons available are frequent with apprehensives: for a language where the apprehensive is restricted to 2$\mathrm{^{nd}}$ person, the researcher will change the grammatical person to 2$\mathrm{^{nd}}$ person when other grammatical persons are suggested in the prompts (except, of course, where restriction on grammatical person itself is the target).

A systematic investigation of the fine-grained semantics of the different functions may uncover subfunctions that have not been clearly identified yet. Users should be aware of the dynamic nature of questionnaires \citep{Lahaussois2019}: this questionnaire is likely to evolve in the next decade, and updates will be posted on the TULQUEST website (\url{http://tulquest.huma-num.fr/en}). 


\section{Apprehensive}
\label{sec:apprehensive}

This section targets the apprehensive function: the marker occurs in main clauses, which are not required to be associated with a preemptive clause (a sentence that gives direction on how to avoid the undesirable event). As mentioned above, one and the same marker might be available for different functions, e.g.\ both for the apprehensive and the precautioning function, the latter being necessarily associated (and often syntactically subordinated) to a preemptive clause (see \sectref{sec:precautioning}). Note that it might be very tricky to determine whether the marker occurs in subordinated clauses and also in main clauses if an insubordination process is in progress. (See \sectref{intro:syntacticstatus} in the introduction to this volume, \citealt{chapters/introduction}.)


\subsection{Intrinsic undesirability}
\label{subsec:apprintrinsic}

For apprehensionals, the expectation is that they are most frequently found in descriptions of events that are universally or culturally judged to be undesirable, which in turn makes warnings a natural context for their occurrence pragmatically (\ref{1.1.1}). The following three questions check for the presence of the undesirability component in the potential candidates for an apprehensive marker. Prompts \REF{p1}--\REF{p3} illustrate occurrences with typically undesirable events, while \REF{p4}--\REF{p5} and \REF{p6}--\REF{p8} illustrate occurrences with typically positive and neutral events, respectively. If a language has no dedicated marker, the questions in (\ref{1.1.1}) will help to find out what is/are the non-dedicated construction/s to express a (highly) potential, undesirable event.

\subsubsection{Is there a construction frequently used in warnings, or in justification/avoidance of an action that invokes a potential negative consequence?}
\label{1.1.1}

\ea \label{p1} {\op}Watch out\textnormal{/}be careful{\cp} you might fall\textnormal{/}slip! \\
\textnormal{Possible context: The addressee is walking very fast on a muddy path.}
\z

\ea \label{p2} Their dog\textnormal{/}a snake might bite me!\\
\textnormal{Possible context: I am not going to their house / I am not going onto that path in the dark.}
\z


\ea \label{p3} He might hit the baby! \\
\textnormal{Possible context: A toddler might hit his younger sibling.}
\z

\subsubsection{Is the reading with \textbf{typically positive situations} like `smile', `laugh', `help', `find a 100 dollar bill' necessarily negative?}
\label{1.1.2}

\ea \label{p4} {\op}Watch out{\cp} you might laugh! \\
\textnormal{Possible context: During an official discourse when one should stay serious and quiet.}\\
\textnormal{Impossible context:\footnote{An impossible context is a context in which the target sentence is expected to be (semantically) infelicitous if the marker has the meaning component tested for, namely undesirability in the present case.} pure potentiality, you might laugh, and that would actually be desirable, or at least not undesirable.}
\z

\ea \label{p5} {\op}Watch out{\cp} Thomas might help! \\
\textnormal{Possible context: Thomas is so clumsy that the speaker would prefer him to not help.} \\
\textnormal{Impossible context: Thomas' helping is a desirable possibility.}
\z


\subsubsection{Is the reading with \textbf{typically neutral situations} like `walk', `talk', `sleep', `come', `cook' necessarily negative?}

This question checks again whether the apprehensional candidate necessarily encodes undesirability.
\label{1.1.3}

\ea \label{p6} {\op}Watch out{\cp} you might fall asleep!\\
\textnormal{Possible context: People being on watch, people watching a movie. }\\
\textnormal{Impossible context: The addressee's sleeping is desirable, e.g.\ if s/he usually has trouble sleeping.}
\z

\ea \label{p7} {\op}Watch out{\cp} Zoe might come! \\
\textnormal{Possible context: The speaker (and addressee) fears Zoe could hear them speaking or see them doing something which would make her angry.} \\
\textnormal{Impossible context: The speaker and/or addressee will rejoice to possibly meet Zoe again.}
\z

\ea \label{p8} {\op}Watch out{\cp} Dan might cook! \\
\textnormal{Possible context: Dan is a bad cook and it is undesirable that he cooks.} \\
\textnormal{Impossible context: Dan is a great cook and the perspective of him possibly cooking is a reason for rejoicing.}
\z

\subsection{Degree of preventability of an event}
\label{subsec:apprpreventability}

This question and the next one test whether the target language restricts the apprehensional morphology to situations where the undesirable event is fully preventable: while one can prevent someone from cooking, one cannot prevent rain or thunder. Some languages allow apprehensives to occur with raining or thundering events since one can still prevent the \textit{consequences} of such an undesirable situation, e.g.\ getting wet (see also \cite[274]{vuillermet2018a}, and \cite{chapters/fedotov}). In Mwotlap (Austronesian; Vanuatu) however, the equivalent of ‘watch out there will be an earthquake' is ungrammatical (p.c.\ A.\ François, May 2014).

Note that there might be a continuum in the preventability of the undesirable event. Nobody can fully prevent the event of falling, but being aware of the risk of falling should lead you to be more careful and potentially control for that risk. Also, the addition of a semantic participant to a verb phrase can facilitate the preventability of the event: while `he might vomit' is an event difficult to prevent per se, `he might vomit \textbf{on you}' is more readily preventable, e.g.\ by moving out of the way.

Note also that the degree of preventability may well be linked to the origin of the marker: if `\textit{beware} you might fall' and `\textit{make sure} you do \textit{not} fall' seem fairly interchangeable, the `beware' formulation focuses on the potential risk \textit{the addressee is not aware of}, while the `make sure not X' formulation focuses on the \textit{actual preventability} of the undesirable event.


\subsubsection{Are non-preventable events compatible with the apprehensive function, e.g.\ events typically expressed by \textbf{impersonal} (like meteorological) verbs and \textbf{[$-$agentive]} verbs?}
\label{1.2.1}

\subsubsubsection{\textbf{Impersonal verbs} like `rain', `snow', `wind.blow', `thunder', `be.an.earth\-quake'; existential verbs}
\label{1.2.1.1}

\ea \label{p9} {\op}Watch out{\cp} it might rain\textnormal{/}thunder!
\z

\ea \label{p10} {\op}Watch out{\cp} there might be snakes on the path!
\z


\subsubsubsection{\textbf{[$-$agentive] verbs}} 
Note that the (un)availability of such verbs might also be linked to the different semantics available with different grammatical persons, see \sectref{subsec:apprgrammaticalpersons}.
\label{1.2.1.2}

\ea \label{p11} {\op}Watch out{\cp} you might fall asleep\textnormal{/}fall\textnormal{/}cough\textnormal{/}vomit!\footnote{Some languages may differentiate between an agentive and non-agentive use of some verbs like `cough', e.g.\ by case or person marking (split S/active-inactive); in that case make sure to test the non-agentive version.}
\z

\ea \label{p12} {\op}Watch out{\cp} Alice might fall asleep\textnormal{/}fall\textnormal{/}cough\textnormal{/}vomit!
\z


\subsection{Availability of the grammatical persons}
\label{subsec:apprgrammaticalpersons}

The availability of the grammatical persons in an apprehensive construction varies very much from one language to the other. In several languages, not all grammatical persons are available. The questions test whether core argument roles in a given construction are available to all or just a subset of grammatical persons.

An extreme case is Matsés (Pano-Tacanan, Peru \& Brazil), where not only the markers do not exist for all grammatical persons, but there are different forms for different person configurations \citep[439ff]{fleck2003}. In such a case, the researcher might want to test all possible configurations.

\subsubsection{What grammatical persons are available with the apprehensive marker?}
\label{1.3.1}

\ea \label{p13} {\op}Watch out{\cp} I\textnormal{/}he\textnormal{/}she\textnormal{/}it\textnormal{/}we\textnormal{/}they might fall on\textnormal{/}hurt you!
\z

\ea \label{p14} {\op}Watch out{\cp} you might fall on\textnormal{/}hurt me\textnormal{/}her\textnormal{/}him\textnormal{/}it\textnormal{/}us\textnormal{/}them!
\z

\ea \label{p15} {\op}Watch out{\cp} you\textnormal{/}he\textnormal{/}she\textnormal{/}it\textnormal{/}we\textnormal{/}they might get lost! \textnormal{or} {\op}OMG{\cp} I might get lost!
\z


\subsection{Interaction of grammatical persons with communicative functions}
\label{subsec:apprinteraction}

While expressing a warning or a general concern is one of the most common functions of the apprehensive, expressing a threat is especially common with 1$\mathrm{^{st}}$ person, e.g.\ \textit{watch out I'll hit you (if you keep bothering me)}, as pointed out by \citet[631]{epps2008} for Hup (Nadahup; Colombia \& Brazil). Interestingly, there is at least one language (Matsés) where the apprehensional expression corresponding to `watch out I'll hit you' necessarily means that the speaker will \textit{accidentally} or \textit{unwittingly} hit the addressee: only the warning or general concern reading is available, while the threat reading is not \citep[439; p.c.\ 2014]{fleck2003}.

In other languages like in Ese Ejja \citep[274]{vuillermet2018a}, the same expression can be used in both warnings and threats. Note that first person subject and full agentivity are typical correlates of the necessary broader conditions that must hold for a threat reading (Searle 1975, Salgueiro 2010), but not the only ones (see the discussion in the Introduction, \sectref{intro:functions}, \citealt{chapters/introduction}).


\subsubsection{What is the reading of [$+$agentive] verbs with 1$\mathrm{^{st}}$ person?}

Note that a warning reading in \REF{p16} is (almost) impossible, while both readings (threat and warning) are available in \REF{p17}, where the hitting event can be [$-$agentive] if non-intentional.
\label{1.4.1}

\ea \label{p16} {\op}Watch out{\cp} I might follow you!
\z

\ea \label{p17} {\op}Watch out{\cp} I might hit you!
\z

\subsubsection{What is the reading of [$-$agentive] verbs with 1$\mathrm{^{st}}$ person?}
If the verb is truly non-agentive (i.e.\ no ``voluntary coughing'' reading available), could it be anything else than a warning?
\label{1.4.2}

\ea \label{p18} {\op}Watch out{\cp} I might fall {\op}over you{\cp}!
\z

\ea \label{p19}(Watch out{\cp} I might cough {\op}on you{\cp}!
\z


\subsection{Semantics: Illocutionary act or expression of general concern}
\label{subsec:apprsemantics}

In some languages, the morpheme might express an undesirable possibility (from the perspective of the speaker), and not necessarily expect a reaction of the addressee. The absence of expectation towards the addressee might be best explained to the consultant by letting them imagine talking to themselves, or comment on a movie where a protagonist is taking risks.

\ea \label{p20} {\op}OMG{\cp} I might lose my temper! \\
\textnormal{Context: The speaker is in an important discussion where they should not lose their temper.}
\z

\ea \label{p21} {\op}OMG{\cp} I might not understand a thing \ldots \\
\textnormal{Possible context: the speaker fears that they might not understand what a teacher or an authority person (e.g.\ a police officer) means.}
\z

\ea \label{p22} {\op}OMG{\cp} Jessica might fall! \\
\textnormal{Possible context: The speaker is suddenly seeing Jessica climbing up a tree, but cannot reach her as they are too far apart. Or the speaker is watching a film where a protagonist is taking risks.}
\z


\subsection{Compatibility with different temporal settings}
\label{subsec:apprtemporal}

Apprehensive markers are usually available in immediately undesirable situations, i.e.\ with future orientation, but in some languages different temporal settings are also available.

\subsubsection{Test the grammaticality with temporal adverbs like \textit{today} vs.\ \textit{tomorrow} vs.\ \textit{yesterday}.}
\label{1.6.1}

\subsubsection{Is the marker available for past events?}
\label{1.6.2}

Apprehensive markers are sometimes available in a past setting, which places them into the domain of epistemic markers: the undesirable event has potentially happened already, but one will not immediately know if it really did or not. In To'abaita (Austronesian; Solomon Islands; \cite[296]{lichtenberk1995}), `be.sick+\textsc{appr}' may mean `you may have been sick' / `you may be sick' / `you may get sick'.

\ea \label{p23} They might have eaten already.
\z

\ea \label{p24} He might have died yesterday.
\z

\ea \label{p25} I might have gotten lost!
\z

\subsubsection{Can the event only be about to happen or also happen some time in the future?}
\label{1.6.3}

Apprehensive markers might only be available with immediate future/prospective orientation but not a general future orientation: in some languages it is possible to warn of an immediate danger, but not of a potential event later in the future. In Hup \citep[609]{epps2008}, an additional future marker emphasizes the possibility of the event to happen in a less immediate future.

\ea \label{p26} {\op}Watch out{\cp} you might hurt yourself! \\
\textnormal{Possible context: if you continue playing with the machete.}
\z

\ea \label{p27} {\op}Watch out{\cp} you might hurt yourself one day! \\
\textnormal{Possible context: if you keep driving when you are drunk as you did last night.}
\z

\subsection{Compatibility with negation}
\label{subsec:apprnegation}

In some languages, the apprehensive marker is not compatible with a negation marker (possibly because of an origin in a negative marker).

\subsubsection{Is the apprehensive construction compatible with overt negation?}
\label{1.7.1}

\ea \label{p28} We might not hear the baby crying!
\z

\ea \label{p29} {\op}OMG{\cp} the luggage might not all fit in the trunk.
\z

\ea \label{p30} {\op}Watch out{\cp} Dan might not come! {\op}\textnormal{although he said so}{\cp}
\z

\ea \label{p31} {\op}Beware{\cp} it might not rain! {\op}\textnormal{when the speaker or addressee wants the rain to water her/his plants}{\cp}
\z

\ea \label{p32} {\op}Watch out{\cp} the dog might not bark at the thieves! {\op}\textnormal{and so not advise when there are thieves}{\cp}
\z

\ea \label{p33} {\op}Watch out{\cp} the cake might not fit into the bag! {\op}\textnormal{and the cake would be ruined if pushed too hard into the bag}{\cp}
\z



\subsection{Compatibility with interrogatives}
\label{subsec:apprinterrogatives}

Many apprehensives are proper mood markers and are impossible in an interrogative sentence.

\subsubsection{Is the apprehensive function also available in a question?}
\label{1.8.1}

\ea \label{p34} Might Bob leave\textnormal{/}laugh\textnormal{/}fall? {\op}\textnormal{and I don't want that}{\cp}
\z

A marker that expresses (among other meanings) undesirability in Kukama-Kukamiría (Tupían; Brazil, Colombia \& Peru) is very frequent in questions \citep[415--417]{Vallejosyopán2016}. This marker also expresses surprise, in the sense that the undesirable event cannot be easily explained: in \REF{p35}, the presence of a new tree that has appeared overnight makes the situation scary. Note that the questions in \REF{p35} and \REF{p36} are not typical questions expecting an answer from the addressee, but rather rhetorical ones expressing incredulity in the face of an event which is so unexpected that it becomes scary.

\ea \label{p35} How might this tree have grown overnight?
\z

\ea \label{p36} How can this bill possibly be in my pocket? \\
\textnormal{Possible context: the speaker has very unexpectedly just found a bill in their pocket and they fear someone could take them as a thief.}
\z



\subsection{Optative uses}
\label{subsec:approptative}

An optative is used to express a wish, regret, hope or desire of the speaker \citep{Grosz2012,dobrushinaetal2013}. A negated optative could then potentially express the opposite of a desire, i.e.\ an undesirable event.

\subsubsection{Does the language have an optative construction, and if so, can it be negated?
If a language displays both an apprehensive and an optative that can be negated, how do the speakers explain the difference between the two forms?} Hypothesis: there is an expectation for the addressee to act against the potential undesirable event in apprehensives, vs.\ absence of expectation for the optative.
\label{1.9.1}


\ea \label{p37} {\op}Watch out{\cp} John might fall! \hspace{0.9cm} vs.\ \hspace{1cm} May John not fall!
\z

\ea \label{p38} {\op}Watch out{\cp} you might fall! \hspace{1cm} vs.\ \hspace{1cm} May you not fall!
\z



\subsection{Scope/distribution}
\label{subsec:apprscope}

Apprehensives might occur with an NP forming a complete utterance which warns of an undesirable event metonymically associated with the tagged undesirable referent. Note that \sectref{sec:timitive} on the timitive function is entirely dedicated to apprehensionals which have scope over NPs.

\subsubsection{Does the apprehensive necessarily attach to a verb or can it also attach to an NP?}
This possibility might depend on the morphosyntactic nature of the morpheme, as in Urarina (isolate; Peru; \cite{Olawsky2006}) where it is a clitic.
\label{1.10.1}

\ea \label{p39} {\op}Watch out{\cp} you'll fall!
\z

\ea \label{p40} {\op}Watch out for{\cp} the crocodile!
\z


\subsection{Compatibility with the evidential system}
\label{subsec:apprevidential}

In Hup \citep[632]{epps2008}, only the reportative is compatible with the apprehensive marker, while no other evidential specifications are.

\subsubsection{Is the apprehensive compatible with the evidential markers of the language (if any)?}
\label{1.11.1}

\ea \label{p41} {\op}I heard{\cp} Rob might die.
\z

\ea \label{p42} {\op}I witnessed{\cp} Rob might die.
\z



\section{Precautioning}
\label{sec:precautioning}

This section targets the precautioning function, in which the apprehensional clause (\textsc{prec}autioning) is obligatorily associated with a clause encoding a preemptive action, which seeks to prevent the (consequences of the) apprehensional clause. In many languages, the precautioning clause is syntactically subordinated to the preemptive clause. The description of the precautioning construction in the target language should indicate the nature of the link between the preemptive and the precautioning clauses:

\begin{itemize}
    
%\begin{compactitem}
    \item At the syntactic level: the precautioning clause is a syntactically subordinate construction if language-internal criteria for subordinate status apply. Alternatively, the link between the two clauses may rather be of a pragmatic nature;
    \item At the semantic level: the precautioning clause is causally linked to the preemptive clause (\citegen{lichtenberk1995} \textsc{preventive} aka.\ \textsc{negative purpose} subtype), or not (aka.\ `\textsc{in-case}' subtype), see \sectref{subsec:prepreventability}. for a detailed explanation of the distinction.
%\end{compactitem}

\end{itemize}

\subsection{Intrinsic undesirability}
\label{subsec:preintrinsic}

The following two questions check for the presence of the undesirability component in the potential candidates for a precautioning marker. \sectref{2.1.1} is about intrinsically undesirable events, and \sectref{2.1.2} and \sectref{2.1.3} about semantically positive and neutral events, respectively. As mentioned above in \sectref{subsec:apprintrinsic}, if a language has no dedicated marker, the first questions in \sectref{2.1.1} might help to find out what construction(s) the target language uses to express the precautioning function.

\subsubsection{Is there a construction used to express both a potential undesirable event and an associated preemptive clause?}
\label{2.1.1}

\ea \label{p43} I quickly picked up the child from the ground {\ob}lest the cow trample on him{\cb}.
\z

\ea \label{p44} Alex should repair the damaged fence {\ob}lest people fall off the cliff{\cb}.
\z

\subsubsection{Is the reading with \textbf{typically positive situations} like `smile', `laugh', `help', `find a 100 dollar bill' necessarily negative?}
\label{2.1.2}

\ea \label{p45} I put my hand on Martina's mouth lest she laugh. \\
\textnormal{Possible context: The protagonists are hiding and making fun of someone, and should not be discovered. }\\
\textnormal{Impossible context: The laughing is desirable.}
\z

\ea \label{p46} They chased Dan away lest he help. \\
\textnormal{Possible context: Dan is so clumsy that they prefer her/him not to help. }\\
\textnormal{Impossible context: The helping is desirable.}
\z


\subsubsection{Is the reading with \textbf{typically neutral situations} like `walk', `talk', `sleep', `come', `cook' necessarily negative?}
\label{2.1.3}

\ea \label{p47} I kept pinching Marine lest she fall asleep. \\
\textnormal{Possible context: An important event where sleeping is not welcome.} \\
\textnormal{Impossible context: The falling asleep is desirable.}
\z

\ea \label{p48}Eva hurried up to leave lest Peter come with her. \\
\textnormal{Possible context: Eva doesn't like Peter \textnormal{/} want Peter to join.} \\
\textnormal{Impossible context: Peter's coming is desirable to Eva.}
\z


\subsection{Degree of preventability of an event}
\label{subsec:prepreventability}

\citet[298]{lichtenberk1995} distinguishes between

\begin{itemize}
    \item the \textsc{preventive} (aka.\  \textsc{negative purpose}, e.g.\ \textit{He gave them his umbrella \textbf{so that} they do \textbf{not} get wet})
    \item and the `\textsc{in-case}' subfunctions (e.g.\ \textit{He gave them his umbrella \textbf{in case} it might rain \textnormal{/} ?\textbf{so that} it does not rain}).
\end{itemize}

\noindent In the first type, there is a causal link between the preemptive event and the non-occurrence (prevention) of the undesirable event encoded by the precautionary clause. In the second type, the undesirable event cannot be prevented itself, but its potentially negative consequences can. According to \citet[308]{lichtenberk1995}, at least one language, Martuthunira (Pama-Nyungan; Australia) distinguishes morphologically between the two subfunctions.

\subsubsection{Must the undesirable event be fully preventable?}
\label{2.2.1}

\subsubsubsection{\textbf{Impersonal verbs} like `rain', `snow', `wind.blow', `thunder', `be.an.earth\-quake'; existential verbs}
\label{2.2.1.1}

\ea \label{p49} Take your coat lest it rain (\#so that it does not rain). \\
$[$\textnormal{vs.} Take your coat lest you get wet (= so as not to get wet).$]$
\z

\ea \label{p50} Check your shoes (before you put them on) lest there be a scorpion inside. (\#so that there will be no scorpion inside.)
\z



\subsection{Who is the judge of the undesirability?}
\label{subsec:prejudge}


An important parameter considered by \citet{green1989} to distinguish between apprehensive and precautioning constructions is the identity of the judge of the undesirable potential event: in apprehensive constructions, only the speaker can judge the event to be undesirable (and the apprehensional construction has thus a modal value), while in precautioning constructions, the subject of the preemptive clause can, or, in other words, ``the speaker simply reports the apprehension of the main clause subject, without himself endorsing or otherwise commenting on the feeling'' (p.\ 168). In precautioning clauses, the undesirability judge is in fact always the subject of the preemptive clause if it is an assertive clause. However, if the preemptive clause is a directive or a modal construction with directive force, then the speaker rather than the subject can be the judge of the potential undesirability, as in \textit{take your coat lest you catch a cold}.

The (non-)availability of the subject to be the judge of undesirability might be considered as a proxy for the distinction between apprehensive and precautioning construction, although this parameter and the obligatory presence of a preemptive clause are theoretically orthogonal. This apparent contradiction comes from the fact that apprehensives tend to be modal (and thus allow only the speaker to be the judge of undesirability) \textit{and} not require a preemptive clause, while, conversely, precautioning markers tend to allow the subject of assertive preemptive clauses to be the judge. The alignment of the two parameters is however not obligatory, so it should not be considered as a test to distinguish apprehensives from precautioning constructions. Once again, the origin of the marker plays a big role in the functions and properties it acquires.

One way to verify the availability of the subject as the judge is to identify the kind of preemptive clause available. If the subject can be the judge, assertive preemptive clauses with third person subject like \REF{p52} should be available. When only the speaker can be the judge,  preemptive clauses are directives as in \REF{p51a}--\REF{p51c}. In \REF{p51a}--\REF{p51c}, the speaker is the one who anticipates the potential falling event, and the apprehensional has some modal value as in apprehensive constructions. In \REF{p52}, Peter, who is the subject of the preemptive clause, is also the one anticipating that she could fall, and the example has no modal value.

\ea \label{p51}
\ea \label{p51a} Tie her shoelaces lest she fall! \\
\ex \label{p51b} You should tie her shoelaces lest she fall!\\
\ex \label{p51c} Don't let her go in socks lest she fall!
\z
\z

\ea \label{p52} Peter tied Natalia's shoelaces lest she fall.
\z




\subsection{Compatibility with negation}
\label{subsec:prenegation}

Probably due to diachronic reasons, the precautioning marker might be incompatible with negation markers.

\ea \label{p53} Keep the door open, otherwise the children will not go out \textnormal{/ ?}lest the children should not go out.
\z

\ea \label{p54} Turn off the music, otherwise we might not hear the dogs barking \textnormal{/} the children calling  \textnormal{/ ?}lest we not hear the dogs barking \textnormal{/} the children calling.
\z

\subsection{Morphosyntax: Subordinating device with apprehensives}
\label{subsec:premorphosyntax}

Some languages make use of a reported speech construction to use the apprehensive marker in the precautioning function, e.g.\ in Tariana (Arawakan; Brazil \& Colombia; \cite[386]{aikhenvald2003}).


\subsubsection{Is the/one of the precautioning construction(s) in the language identical to an apprehensive (\sectref{sec:apprehensive}) with a quotation construction (see the equivalent literal translation in square brackets below the examples)?}
\label{2.5.1}

\ea \label{p55} You covered my mouth with your hand, lest they hear me.\\
\textnormal{$[$equivalent literal translation: You covered my mouth, \textbf{thinking/saying}: ``they might hear you!''$]$}
\z

\ea \label{p56} You ran away lest he see you. \\
\textnormal{$[$equ.\ lit.\ tr.: You ran away, \textbf{thinking/saying}: ``he might see me!''$]$}
\z

\ea \label{p57} They don't eat sugar lest they get bad teeth. \\
\textnormal{$[$equ.\ lit.\ tr.: They don't eat sugar, \textbf{thinking/saying}: ``we might get bad teeth!''$]$}
\z


\subsection{`Fear' complementation function}
\label{subsec:prefear}

\citet{lichtenberk1995} has observed that apprehensionals available in the precautioning function could also be available for what he calls the `fear' function: apprehensional complementizers which appear with clauses that serve as the non-subject complement of a transitive predicate of fear or concern. For instance, (archaic) English \textit{lest} may mark subordinate clauses that complement main clause predicates expressing events of concern like `be frightened', `be anxious', `be worried', etc., as in \textit{I am frightened} \textsc{lest} \textit{he come}. Expectedly, its combination with main predicates expressing positive emotional states is impossible, e.g.\ \textit{*I am happy} \textsc{lest} \textit{he come}. Note that most languages have a general complementizer (e.g.\ \textit{that} in English) available with main predicates expressing fear or concern, as well as any other emotional states (and often other predicate types expressing for instance cognition or perception). English has thus both a general complementizer as in \textit{(I'm happy/I'm worried/I heard)} \textsc{that} \textit{he will/might come}, in addition to its apprehensional marker \textit{lest} available with `fear' or `concern' predicates only.

The precautioning and the `fear' complementation functions are unified by the semantics of undesirability, and by their subordinated syntax. While there exist three logically possible constellations of available functions for an apprehensional construction, only the first two are attested so far:

\begin{itemize}
    \item both the precautioning and the `fear' complementation functions;
    \item the precautioning function only;
    \item the `fear' complementation function only.
\end{itemize}

\noindent A language might allow both the precautioning and the `fear' complementation function despite the redundancy of the information on the undesirability of the event in the latter function. This is actually precisely the semantic harmony which makes the collocation of a `fear' predicate and an apprehensional clause so common in languages (see e.g.\ \cite{López-Couso2007} on English \textit{lest}). By contrast, having an apprehensional marker dedicated to complementizer function with predicates of fear or concern seems non-economical, as a general subordinate marker could fulfill that function, especially if no precautioning marker is available in the language.

It is still worth checking whether unexpected diachronic reasons could have given rise to an apprehensional marker that would only be used in this `fear' function. The questions below examine the availability of the `fear' function to the precautioning marker.


\subsubsection{Does the precautioning marker -- or, very unexpectedly, a marker dedicated to `fear' complementation -- exclusively introduce clauses that serve as the non-subject complement of a predicate of fear or concern?}
\label{2.6.1}

\ea \label{p58} I am afraid that\textnormal{/}lest he come.
\z

\ea \label{p59} I am happy that\textnormal{/*}lest he come(s).
\z

\ea \label{p60} I heard that\textnormal{/*}lest he come(s).
\z

\noindent Here the researcher will have to use language-internal criteria to determine whether the apprehensional clause is a syntactic complement of the predicate of fear or concern. For instance, in \textit{I'm afraid, he might come} (with an intonational break), the apprehensional clause is not a syntactic complement, as opposed to \textit{I'm afraid, he might come} (with no intonational break), where it is.

\subsubsection{Are events of prevention (expressed by verbs like `forbid', `impede') also available with the precautioning (or `fear' complementation) marker?}
\label{2.6.2}

\ea \label{p63odd} I forbid you to go there {\op}\sim I forbid you \textsc{lest} you go there{\cp}.\textnormal{\footnote{In Mwotlap \citep[304--305]{francois2003}, such a sentence is possible, while it is not in English.}}
\z




\section{Timitive}
\label{sec:timitive}

This section targets the timitive function, i.e.\ when an apprehensional marker is associated with an NP (and possibly also with a VP, see \sectref{subsec:timmulti}). Here, the undesirability is about an entity, and its anticipated dangerous or at least unwelcome behavior. 



\subsection{Intrinsic undesirability}
\label{subsec:timintrinsic}

Undesirable entities may be marked by a general causal marker as in \textit{He ran away} \textbf{because of} \textit{the bees}, or by a dedicated causal marker for undesirability, here referred to as timitive, as in \textit{He ran away} \textbf{for fear of} \textit{the bees}, with a grammaticalized construction corresponding to English \textit{for fear of} (cf. example \ref{ex:v1c}). The following two questions thus check for the contexts where the timitive marker may appear: in \sectref{3.1.1}, with an intrinsically undesirable entity, and in \sectref{3.1.2} with typically positively evaluated entities. Prompt \REF{p61} might help to find out what the (non)-dedicated construction(s) are that express that an undesirable/feared entity is the cause or justification for an action.

\subsubsection{Is there a construction frequently used in contexts where an undesirable entity is mentioned?}
\label{3.1.1}

\ea \label{p61} I ran away \textbf{for fear of/because of} the fierce dog.
\z

\subsubsection{Is the reading with a typically positively evaluated entity like `mother', `husband' or `friend' necessarily negative?}
\label{3.1.2}

\ea \label{p62} I cooked \textbf{for fear of/because of}mother\textnormal{/}husband. \\
\textnormal{Possible context: my mother/husband will hit me if I have not cooked food for her/him.} \\
    \textnormal{Impossible context: It is my mother/husband's birthday and I am happy to cook for her/him.}
\z

\ea \label{p63} He ran away \textbf{for fear of/because of} his friend.\\
\textnormal{Possible context: the friend has become aggressive.}\\
\textnormal{Impossible context: the friend has always been a good friend and it has not changed.}
\z




\subsection{Verb types}
\label{subsec:timverbtypes}

\subsubsection{Semantic role of the entity: Entity to be avoided or stimulus of fear}
\label{3.2.1}

This question and the next one test the detailed semantics of the timitive. The marker might encode \textit{avoidance} or \textit{fear} of an entity. In both cases, the entity is undesirable, but each has its own very distinct role:

\begin{itemize}
    \item in the case of \textit{avoidance}, its role is to be the avoided entity. As a consequence, the main event is a preemptive action to avoid that undesirable entity, as in \textit{I ran away \textbf{for fear of} the snake/\textbf{to avoid} the snake};
    \item in the case of \textit{fear}, by contrast, its role is to be a stimulus for a behaviour or reaction. This is why the main event is not necessarily construed as a preemptive action, and might for instance be a physiological reaction triggered by fear, as in \textit{I fainted \textbf{for fear of} the snake/\textbf{*to avoid} the snake}.
\end{itemize}


That semantic distinction can thus be tested with verbs of physiological reactions indicative of fear but not conducive to avoiding the entity. In Marrithiyel, the timitive seems to require a preemptive action \citep[58]{green1989}, unlike in Ese Ejja, where a physiological reaction is also available \citep[284]{vuillermet2018a}. The prompts in \REF{p66} specifically test whether it can also encode disgust, since the emotion of disgust is close to fear.\footnote{Fear and disgust are both human reactions triggered by a threatening situation or entity.}

\ea \label{p64}
\ea \label{p64a} He ran away \textbf{for fear of}/\textbf{to avoid} the bees. \\
\ex \label{p64b} He hid \textbf{for fear of}/\textbf{to avoid} the bear.
\z
\z

\ea \label{p65}
\ea \label{p65a} He burst into tears \textbf{for fear of}/\textbf{*to avoid} the bees \\
\ex \label{p65b} He fainted \textbf{for fear of}/\textbf{*to avoid} the bear.
\z
\z

\ea \label{p66}
\ea \label{p66a} I went away \textbf{out of disgust for}/?\textbf{for fear of} the vomit. \\
\ex \label{p66b} I vomited \textbf{out of disgust for}/?\textbf{for fear of} the stinky fish.
\z
\z

\subsection{Possible entities}
\label{subsec:timentities}

Up to now the literature has not reported restrictions on the nature of the undesirable entity, but it is not excluded that a timitive in a given language could be unavailable with certain types of entities. The following questions target the availability of (in)animate and (non-)human entities.

\ea \label{p67} Rebecca left \textbf{for fear of} her mother-in-law.
\z

\ea \label{p68} He crossed the bridge \textbf{for fear of} the plants {\op}\textnormal{e.g.} nettle{\cp}.
\z

\ea \label{p69} I went this way \textbf{for fear of} the thorny scrub\textnormal{/}the excrement.
\z

\ea \label{p70} He stayed in \textbf{for fear of} the rain\textnormal{/}thunder\textnormal{/}snow\textnormal{/}wind\textnormal{/}cold.
\z

\subsection{Possible TAM situations}
\label{subsec:timTAM}

Up to now the literature has not reported restrictions on the TAM values of the main event, but it is not excluded that a timitive in a given language could be unavailable with some TAM markers, which is what the following prompts test for. (The previous prompts are all in the past to test for only one parameter at a time.)

\ea \label{p71} I am preparing the dinner \textbf{for fear of} my husband/mother.
\z

\ea \label{p72} He will come back quickly \textbf{for fear of} the night.
\z

\subsection{Compatibility with interrogative words}
\label{subsec:timinterrogative}

There should be no restrictions on whether interrogative words are available with any NP markers, but it cannot be excluded; this is what the following prompt tests for.

\ea \label{p73} \textbf{For fear of} what did you run away?
\z


\subsection{Multifunctional apprehensionals}
\label{subsec:timmulti}

The last two questions are potentially redundant with \sectref{subsec:apprscope} above, but one can imagine a language with more than one apprehensional, where the various functions of each apprehensional might overlap. Note that \sectref{3.6.2} might be the necessary step for a timitive to acquire the apprehensive function.

\subsubsection{Can the timitive attach to NPs only or can it also attach to nominalized/subordinated verbs/VPs/clauses?}
\label{3.6.1}

\ea \label{p74} He screamed \textbf{for fear of} falling \textnormal{/} [the baby falling].
\z

\ea \label{p75} He ran away \textbf{for fear of} [being eaten\textnormal{/}attacked\textnormal{/}harmed.]
\z

\subsubsection{Can the timitive NP form a non-verbal clause? \normalfont{(cf.\ in Manambu [mana1298] (Ndu; Papua New Guinea), \cite[153--154]{aikhenvald2008manambu})}}
\label{3.6.2}

Note that such non-verbal clauses are similar to apprehensive clauses in their illocutionary reading (warning or threat).

\ea \label{p76} \textbf{Watch out for} the tiger\textnormal{/}the excrement!
\z

\ea \label{p77} \textbf{Beware of} my strength!
\z


\section*{Abbreviations}
\begin{tabularx}{.5\textwidth}{@{}lQ@{}}
\textsc{3a} & 3$\mathrm{^{rd}}$ agent \\
\textsc{abs} & absolutive \\
\textsc{appr} & apprehensive \\
\textsc{apprl} & (general) apprehensional marker \\
\textsc{erg} & ergative \\
\textsc{fut} & future \\
\textsc{irr} & irrealis \\
\textsc{mid} & middle \\
\textsc{neg} & negation \\
\textsc{nirr} & non-irrealis \\
\textsc{nsg} & non-singular \end{tabularx}%
\begin{tabularx}{.5\textwidth}{@{}lQ@{}}
\textsc{o} & object \\
\textsc{OMG} & oh my God \\
\textsc{poss} & possessive \\
\textsc{prec} & precautioning \\
\textsc{pres} & present \\
\textsc{rdp} & reduplication \\
\textsc{real} & realis \\
\textsc{s} & subject \\
\textsc{sbj} & subject \\
\textsc{sg} & singular \\
\textsc{tim} & timitive \\\\
\end{tabularx}


\section*{Acknowledgements}
Earlier versions of this questionnaire have been worked on during my previous postdoctoral fellowship from LABEX ASLAN (ANR-10-LABX-0081) of the Université de Lyon with the program ``Investissements d'Avenir'' (ANR-11-IDEX-0007). I am very thankful  to Anetta Kopecka and Maïa Ponsonnet for their many insightful comments on an earlier version of the questionnaire, and to Eva Schultze-Berndt and Martina Faller for their comments on the current version. All remaining mistakes are, of course, my own.

\printbibliography[heading=subbibliography,notkeyword=this]
\end{document}
